\documentclass[10pt,a4paper]{amsart}
\usepackage[margin=19mm]{geometry}
\usepackage{amsmath,amssymb,mathtools}
\usepackage[scr=rsfs,cal=euler]{mathalfa}
\usepackage{xfrac}
\usepackage[unicode,hidelinks,bookmarks=false]{hyperref}
\newcommand{\ie}{i.e.\ }
\newcommand{\cf}{cf.\ }
\newcommand{\resp}{respectively\ }
\newcommand{\Subsection}{\subsection}
\providecommand{\nobreakdash}{\nobreak\hbox{-}}
\setlength{\emergencystretch}{2em}
\multlinegap0pt
\usepackage{amssymb}
\usepackage{float}
\newtheorem{lem}{Lemma}
\newcommand{\Ll}{\mathcal L}
\DeclareMathOperator{\mul}{mul}
\DeclareMathOperator{\ran}{ran}
\title{Hermitian Pencils and their Representation in Krein Spaces}
\author{Rabeb Aydi, Omaima Kchaou, Carsten Trunk}
\date{Source: arXiv:2607.05852v1 (CC BY 4.0)}
\begin{document}
\maketitle

\noindent\emph{Source and licence.} Hermitian Pencils and their Representation in Krein Spaces. Version arXiv:2607.05852v1 (CC BY 4.0), distributed by arXiv under the Creative Commons Attribution 4.0 International licence, http://creativecommons.org/licenses/by/4.0/, as recorded in the arXiv licence metadata for this exact version and shown on the abstract page https://arxiv.org/abs/2607.05852v1. Full licence notice: \texttt{licenses/arxiv-2607.05852-CC-BY-4.0.txt}.\par\smallskip

\noindent\textit{Editorial scope.} Counted unit: the article's lem representation (source0.tex lines 314--329). The article's complete original proof of it is delivered with it (source0.tex lines 332--367, one \texttt{proof} environment of 36 lines). No mathematics is written, repaired or re-derived: the statement and the proof are the article's own lines, in the article's own notation, and no retained line is substituted or re-flowed.  Retained uncounted as the source-local context the proof needs: lines 157--163; lines 286--288.  Nothing was omitted from any retained span, so the package carries no omission record.  No retained line contains a citation, so the wrapper carries no bibliography and no bibliography entry.  No retained line contains a figure environment, an image-inclusion command or drawing code, so no image is delivered and the package declares no asset dependency.  Editorial formatting only, in addition: an AMS \texttt{amsart} preamble that restates the article's own declarations and macros used by the retained lines, namely the environments \texttt{lem} on separate counters, together with the standard packages those lines need.  The retained environments print in this document as Lemma 1 in order of appearance, whereas the article prints the same items with the numbering of its own full document.  The complete native source of the article as distributed in the arXiv e-print for this exact version is bundled unchanged as \texttt{sources/204577/source0.tex}.
\begin{equation} \label{range-kernel}
	\mathcal{L} = \ran\begin{bmatrix}
			E\\ A
			\end{bmatrix}, % = \text{Ker} \begin{bmatrix}
			%G, F
			%\end{bmatrix}.
\end{equation} 

\begin{equation} \label{Krein Adjoint of L} 
\mathcal{L}^{+} = \{\{x,y\}\in X \times X~|~ (JA)^{*}x=(JE)^{*}y \}.
\end{equation} 

\begin{lem}  \label{Range representation}
 Let $\mathcal{L}$  be a linear relation as defined in \eqref{range-kernel}. Then 
 \begin{align}\label{Equation 6}
       AE^{-1}= \mathcal{L}= \textup{ran}\begin{bmatrix}
			EJ\\ AJ
			\end{bmatrix}= \textup{ran}\begin{bmatrix}
			EJ\\ -I+\lambda EJ
			\end{bmatrix},
    \end{align} 
 Further, the following holds.
 \begin{enumerate}
    \item $\mathcal{L}$ is closed in $X\times X.$
    \item $\ker\Ll^+= (\ran JA) ^\perp$ and $\mul\Ll^+= (\ran JE) ^\perp.$
    \item  $\ker\Ll= (\ran JAJ) ^\perp$ and $\mul\Ll= (\ran JEJ) ^\perp.$
\end{enumerate} 
\end{lem}

\begin{proof}
 We first prove $(\ref{Equation 6})$. By definition, we have that
 \begin{eqnarray}  \nonumber
 AE^{-1}&=& \{\{x,z\} \in X \times X ~|~ \{x,y\}\in E^{-1}, ~ \{y,z\} \in A \} \\[1ex]\label{Sylvester}
 &=& \{\{x,z\} \in X \times X ~|~  x=E y, ~ z= Ay, \ y \in X \} \\[1ex]\nonumber
 &=&  \textup{ran}\begin{bmatrix}
			E\\ A
			\end{bmatrix} = \mathcal{L}.
 \end{eqnarray}
 A substitution $x=(\lambda E-A)^{-1}y,$ $y\in X$, gives
  \begin{eqnarray*}\mathcal{L}&=& \{\{E(\lambda E-A)^{-1}y,A(\lambda E-A)^{-1}y\} ~|~  y\in X\} .
 \end{eqnarray*} 
 Obviously, one has $AJ=(A-\lambda E+\lambda E)J=-I+\lambda EJ$, so that  
 $$
        \mathcal{L}=  \textup{ran}\begin{bmatrix}
			EJ\\ -I+\lambda EJ
			\end{bmatrix}.
   $$ 
We now prove 1. Let 
$$
\{x_{n},z_{n}\}\in \textup{ran}
\begin{bmatrix}
			EJ\\ AJ
			\end{bmatrix}
$$
such that $x_{n}\rightarrow x$, $z_{n}\rightarrow z$, $n\rightarrow \infty$. It follows from \eqref{Sylvester} that $\{x_{n},z_{n}\}\in AJ(EJ)^{-1}$ and hence there exists $y_{n}\in X$ such that 
$$
\{x_{n},y_{n}\}\in (EJ)^{-1} \ \text{and}\ \{y_{n},z_{n}\}\in AJ.
$$
Then it is clear that 
$$
x_{n}=E Jy_{n}\quad  \mbox{and}\quad  z_{n}=A Jy_{n}=-y_{n}+\lambda EJy_{n}.
$$
We deduce that $y_{n}\rightarrow y$, $n\rightarrow \infty$, for some $y\in X$ with $y=\lambda x-z$.
The boundedness of $EJ$, implies  $\{y,EJy\}\in EJ,$ hence with $x=EJy$, we obtain $\{x,y\}\in(EJ)^{-1}$. From the boundedness of $AJ$ and $\{y_{n},z_{n}\}\in AJ $, we have that $\{y,z\}\in AJ.$ Overall, we conclude $\{x,z\}\in AJ(EJ)^{-1}=\mathcal{L}$. We continue with the proof of 2. From \eqref{Krein Adjoint of L} we obtain $\mul\Ll^+= \ker (JE)^*$ and $\ker\Ll^+= \ker (JA)^* $. Thus,  $\mul\Ll^+= (\ran JE) ^\perp$ and $\ker\Ll^+= (\ran JA) ^\perp$. The last claim follows in a similar way as in $2$.
\end{proof}

\end{document}
